\section{Study Design}
\label{sec:study-design}

\subsection{Outcome-blind corpus construction}
\label{sec:corpus-construction}

The study used a frozen corpus of 39 repository--operator cases
drawn from machine-learning research software. Corpus construction was
performed separately from mutation execution and was intentionally
outcome-blind: repository eligibility, mutation candidates, revisions, and
validation workflows were selected without executing candidate repositories
or observing mutation outcomes.

The corpus consisted of two batches. B01 was a 10-repository calibration batch
used during development of the empirical protocol. B02 was the primary
operator-specific corpus and contained 29 repository--operator cases. The
final B02 composition was 10 \texttt{random-seed}, 10
\texttt{dependency-pin}, 6 \texttt{data-split}, and 3
\texttt{cv-fold-count} cases.

For B02, repositories were screened using an operator-specific static
eligibility procedure. Screening could use repository source code,
documentation, tests, continuous-integration configuration, experiment or
example source, and publication metadata, but could not execute repository
code, install dependencies, create environments, download datasets for
validation, or inspect mutation outcomes. In particular, repositories were
not selected because their tests appeared likely to kill a mutation or
because a candidate appeared especially sensitive to the proposed change.

The stopping rule was fixed before primary execution as

\[
n_{\mathrm{primary}} = \min(10, n_{\mathrm{eligible}}).
\]

Consequently, \texttt{random-seed} and \texttt{dependency-pin} each
contributed ten B02 cases, whereas the eligible static frames for
\texttt{data-split} and \texttt{cv-fold-count} were exhausted at six and
three cases, respectively. These smaller operator frames therefore represent
the eligible population identified under the frozen screening procedure
rather than post-hoc reductions based on execution outcomes.

The complete 39-case corpus, including repository revisions, candidate
mutations, and selected validation workflows, was frozen before primary
mutation execution. No repository was subsequently added to compensate for
setup failures, workflow failures, survived mutations, or other empirical
outcomes.


\subsection{Mutation operators and candidate selection}
\label{sec:mutation-operators}

The study used four reproducibility-relevant mutation operators implemented by
MLReproMutate. Each operator represents a controlled change to an experimental
or environment choice that can affect the reproducibility of an ML workflow.

\texttt{random-seed} targets supported literal calls to Python, NumPy, or
PyTorch random-seed APIs and changes an integer seed from \(N\) to \(N+1\).

\texttt{dependency-pin} relaxes an exact dependency pin. It replaces
\texttt{package}\allowbreak\texttt{==}\allowbreak\texttt{version} with
\texttt{package}\allowbreak\texttt{>=}\allowbreak\texttt{version}.

\texttt{data-split} removes a non-null \texttt{stratify} argument from a
supported \texttt{train\_test\_split} call.

\texttt{cv-fold-count} changes a literal \texttt{n\_splits=N} argument in a
supported KFold-family call to \texttt{n\_splits=N+1}.

Candidate selection was restricted to syntax supported by the frozen operator
definitions. The study therefore does not attempt to enumerate every
reproducibility-relevant choice that can occur in ML research software.
Instead, it evaluates a deliberately narrow set of mechanically identifiable
and auditable mutations.

For each selected case, the repository revision and candidate mutation were
fixed before execution. Candidate selection was not changed after observing
whether a baseline executed successfully or whether a mutant was detected.


\subsection{Frozen validation workflows}
\label{sec:validation-workflows}

Each case was associated with one validation workflow selected before mutation
execution. The objective was not to construct a new test suite for the
repository, but to measure the sensitivity of validation behavior that already
existed in the research artifact.

The protocol allowed five workflow categories:
\texttt{upstream-test}, \texttt{ci},
\texttt{documented-validation}, \texttt{documented-experiment}, and
\texttt{documented-example}. The selected command and its repository
documentation or source reference were frozen together with the case.

These categories describe the provenance or form of the workflow and are
treated categorically rather than ordinally. For example, an upstream test
command is not assumed a priori to provide a stronger reproducibility oracle
than a documented experiment, and an example workflow is not assumed to
provide a weaker one.

The term \emph{validation workflow} is intentionally broader than
\emph{reproducibility safeguard}. A workflow may execute substantial research
code while providing no explicit check on its outputs. Accordingly,
successful completion alone was not interpreted as evidence that a workflow
constituted a strong reproducibility oracle.


\subsection{Validation-oracle classification}
\label{sec:oracle-classification}

To distinguish workflow form from the mechanism that determines workflow
success, protocol version 2 prospectively classified the validation oracle for
every selected B02 case. Oracle classification was recorded before mutation
execution and was therefore independent of \textsc{Killed},
\textsc{Survived}, or \textsc{Equivalent} outcomes.

Four oracle categories were permitted. \texttt{assertion} indicates that
explicit assertions or test expectations determine success.
\texttt{metric-threshold} indicates an explicit numerical acceptance
criterion. \texttt{reference-comparison} indicates comparison with an expected
value or reference artifact. \texttt{completion-only} indicates that the
selected workflow is considered successful if it terminates with exit status
zero, without an explicit assertion, threshold, or reference comparison
relevant to the workflow result.

The frozen B02 corpus contained 22 \texttt{completion-only} and 7
\texttt{assertion} workflows; no B02 case used a
\texttt{metric-threshold} or \texttt{reference-comparison} oracle.

For a secondary descriptive contrast fixed before outcome analysis, the latter
three explicit-oracle categories were grouped as
\texttt{substantive-oracle}, while \texttt{completion-only} remained
separate. This binary grouping does not imply an ordering among
\texttt{assertion}, \texttt{metric-threshold}, and
\texttt{reference-comparison}.

B01 predates the protocol-v2 oracle field. We therefore retain the frozen B01
workflow categories but do not retrospectively treat B01 cases as if an oracle
category had been prospectively recorded. The primary oracle analysis for RQ2
is consequently restricted to B02.


\subsection{Primary execution protocol}
\label{sec:primary-execution}

Primary execution followed a baseline-first protocol. For each frozen case,
the selected repository revision was prepared according to the study protocol
and the frozen validation workflow was first executed on the unmodified
baseline. Mutation execution proceeded only when the baseline satisfied the
protocol's evaluability requirements.

If the baseline could not be established because of environment setup failure,
unavailable workflow requirements, or other infrastructure barriers, the case
was classified as non-evaluable at the primary layer rather than assigned a
mutation outcome. Mutation outcomes were therefore defined only for cases in
which the corresponding baseline and mutant could be evaluated under the
frozen workflow.

A mutation was classified as \textsc{Killed} when the selected validation
workflow failed or returned a non-zero status after the mutation. It was
classified as \textsc{Survived} when the same workflow completed successfully
despite the mutation. A mutation was classified as \textsc{Equivalent} when
semantic verification established that the applied mutation had not changed
the relevant semantics of the evaluated case.

Setup failures, workflow-unavailability events, and other infrastructure
failures were not counted as mutation outcomes.


\subsection{Semantic verification}
\label{sec:semantic-verification}

Mutation survival alone does not establish that a mutation changed the
semantics of the evaluated computation. The analysis therefore separates
validation-workflow detection from semantic equivalence.

Semantic verification distinguished four states: confirmed
non-equivalent, confirmed equivalent, unverified, and not run. Confirmed
equivalent mutations were excluded from the meaningful detection denominator.

The primary detection measure used for RQ1 and RQ2 is

\[
\frac{\mathrm{KILLED}}
     {\mathrm{KILLED} +
      \mathrm{confirmed\text{-}non\text{-}equivalent\ SURVIVED}}.
\]

This distinction is important for interpreting survival. A confirmed
non-equivalent \textsc{Survived} mutation means that the selected frozen
validation workflow did not detect that particular reproducibility-relevant
change. It does not imply that the repository as a whole is irreproducible.


\subsection{Bounded restoration of non-evaluable cases}
\label{sec:d027-restoration}

Primary execution revealed substantial environment and dependency drift in
historical ML research repositories. To distinguish mutation sensitivity from
failures caused by contemporary execution environments, we applied a separate
bounded restoration protocol, D027, to a predefined subset of initially
non-evaluable cases.

D027 was an additional evidence layer rather than a revision of the primary
study. Canonical primary outcomes and failure classifications were retained
unchanged. Restoration preserved the frozen repository revision, mutation
candidate, validation workflow, and oracle classification.

Restoration followed several constraints intended to prevent outcome-directed
repair. Source-level compatibility patches were prohibited; validation
workflows could not be skipped or weakened; the baseline had to be restored
before mutation evaluation; the restoration recipe was fixed before executing
the mutant; and baseline and mutant were evaluated using symmetric restoration
conditions. Restoration attempts were bounded, and infrastructure failures
were kept distinct from mutation outcomes.

The D027 cohort contained 24 cases. A structured D027 report or assessment was
available for 19 cases. Substantive restoration was attempted for 18 cases, 11
of which were successfully restored and subsequently evaluated. Seven cases
were not restored after substantive attempts. One additional case had a D027
report but no substantive restoration attempt because repository identity
recovery failed, and five cohort cases had no D027 report.

Successful D027 evaluations were combined with primary evaluations only at the
analysis layer. D027 never changed the frozen workflow or oracle
classification and never replaced the canonical primary record.


\subsection{Analysis procedure and denominators}
\label{sec:analysis-denominators}

We distinguish three levels of accounting: the complete frozen corpus, the set
of evaluable cases, and the set of confirmed non-equivalent mutations that
enter the meaningful detection denominator.

The complete corpus contains 39 cases. Primary execution evaluated 13 cases
and left 26 non-evaluable. After incorporating successful D027 restoration, 24
cases were evaluable and 15 remained non-evaluable.

Among the 24 combined evaluated cases, one mutation was confirmed equivalent.
The meaningful detection denominator therefore contains 23 confirmed
non-equivalent mutations: 2 \textsc{Killed} and 21 \textsc{Survived}.

RQ1 reports detection over these 23 cases. RQ2 first stratifies the same
combined meaningful cases by the frozen workflow category. The primary oracle
analysis is restricted to B02 because only B02 contains prospectively recorded
oracle classifications. Fifteen B02 cases entered the confirmed
non-equivalent meaningful denominator for this analysis: 2 \textsc{Killed}
and 13 \textsc{Survived}.

RQ2 is treated as descriptive and exploratory. Category sizes are small and
uneven, and only two confirmed non-equivalent mutations were detected. We
therefore report counts and descriptive proportions rather than interpreting
category differences as causal effects or evidence of statistical superiority.

As an additional diagnostic for RQ2 interpretation, we examined operator
composition within the prospectively classified B02 oracle groups and
evaluability attrition across categories. These diagnostics were performed
after the workflow and oracle classification frame had been independently
frozen. They are used to identify potential confounding and denominator
imbalance, not to redefine the RQ2 categories.
